% =====================================================================
% Section 4.4 --- The welfare measure (DRAFT, 2026-10-15)
% ---------------------------------------------------------------------
% Proposal file --- NOT wired into main.tex. Intended as a short new
% subsection of Section 4 (kernel), placed after the notation glossaries.
% If Section 4.3 is moved out to the closure appendix (App. G/B), this
% becomes 4.3.
%
% Job: (i) state the two aggregate measures the paper reports -- real GDP
%      and the welfare index -- and why they are not the same object;
%      (ii) name the alternatives and recommend one.
% Numbers cited are from Table 4 (matrix_5x3_v10): F2 mobile  dY = 0.00 %,
%      dW = -1.82 %; fixed-wage F3  dY = +1.78 %, dW = +2.26 %.
% =====================================================================
\subsection{The welfare measure}\label{subsec:welfare}

The results below report two aggregate measures that must not be conflated.

\emph{Real GDP} is an income-side quantity index: the wage bill
$w\sum_i L_i$ deflated by a seven-component T\"ornqvist price index of final
demand (consumption, government purchases and the programme, investment,
exports, and the three import and tax leaks booked in
Appendix~\ref{app:kernel}). Under a demand-only shock every price is pinned
at its baseline value, so the deflator is exactly one and real GDP moves
with the employment index $w\sum_i L_i$ alone --- which is why the
full-employment mobile rows read a flat $0.000\%$ by construction. Real GDP
measures how much is produced, not how well the household fares.

\emph{The welfare index} is a T\"ornqvist quantity index of the household
consumption bundle against its baseline. It is the constant-returns
final-demand aggregator of the Baqaee--Farhi kernel, and in the closed
economy it coincides with real GDP; in our open economy it does not.
The F2 mobile row is the clean illustration: real GDP is exactly
unchanged, yet the tax-financed programme costs the household
$1.82\%$ of consumption. Reporting both is the point --- a measure that
would read ``GDP-neutral'' and stop there would hide the loss.

Because the representative household has homothetic CES preferences, the
model also delivers an \emph{exact} money-metric welfare change: with the
expenditure function $e(p,u)=u\,\Pi(p)$, the change in real consumption
expenditure $E^h/\Pi(p)$ is the consumption-equivalent (equivalent
variation) of the shock, to all orders. Under the consumption-price
numeraire $\Pi(p)=1$, so this measure is simply the change in after-tax,
post-saving household resources $E^h$. Where prices do not move it agrees
with the reported quantity index exactly (equal proportional expansion of
every $c_i$), and it differs from it only in the price-moving sectoral and
no-reallocation rows, where the T\"ornqvist quantity index is a
second-order approximation to it. One caveat of scope: under F1 the
preference weights $d_i$ are themselves shifted, so an F1 ``welfare'' number
is a statement about the composition of the bundle rather than a
utility comparison across unchanged preferences, and should be read as
such.
A second caveat travels with both: the index prices the consumption bundle
only and is silent on the disutility of the additional work. In the
fixed-wage rows, where employment absorbs the shock (F3: consumption
$+2.26\%$, employment $+1.78\%$), a welfare measure that counted the
leisure margin would report a smaller gain; the elastic-supply closures are
reduced forms without a labour--leisure income effect
(Appendix~\ref{app:labour-closures}), so a leisure-augmented measure would
require an explicit preference over leisure.

\emph{Alternatives.} Three candidates are worth naming. (i) The exact CES
money-metric above, $E^h/\Pi(p)$ --- free in the current kernel, exact, and
identical to the reported index on the demand-only price-pinned rows; we
recommend reporting it as the headline welfare number, with the quantity
index kept as a cross-check. (ii) A leisure-augmented equivalent variation,
which would require specifying a labour--leisure preference and is the
natural home of any future extension that makes the extensive margin
central. (iii) A broader final-demand index that aggregates consumption
\emph{and} the programme ($c+g$), which avoids reading a public-investment
programme through household consumption alone; we would report it as a
supplementary series rather than as the welfare measure, since it embeds a
valuation of public goods the model does not otherwise take a stance on.
